\section{结论}\label{sec:conclusion}

本文综述了朗兰兹纲领从1967年信件到2024年几何朗兰兹证明的发展历程。可以概括出三条主线：

\begin{enumerate}[label=(\arabic*)]
    \item \textbf{互反律主线}：从阿廷 $L$ 函数到 Jacquet--Langlands、Langlands--Tunnell、模性定理与 Serre 猜想，二维互反律在可解与“满约化”意义下已近完成；函数域上 $\operatorname{GL}_n$ 的完全解决（Drinfeld、L. Lafforgue）与 V. Lafforgue 的一般群构造，证明了纲领整体并非幻觉。
    \item \textbf{函子性主线}：基变换、对称幂与秩积的部分结果持续兑现为解析数论红利；迹公式与内窥理论（Ng\^o 基本引理、Arthur 分类）从谱一侧逼近函子性，Shimura 簇与势自守性则从算术一侧逼近——两条路线在 $R=T$ 处会师。
    \item \textbf{几何主线}：Shtuka、Hitchin 丛、完美胚到导出代数几何，几何朗兰兹从函数域试验场成长为独立的数学分支；Fargues--Scholze 的局部几何化与 Gaitsgory--Raskin 的整体证明确立了“范畴化”纲领的严格地位，也为数域情形的算术化指明方向。
\end{enumerate}

朗兰兹纲领是最能体现“数学深刻统一性”的活体标本：一条非阿贝尔互反律的猜想，最终把数论、表示论、代数几何与量子场论焊接成同一张 $L$ 函数之网。真二十面体 Artin 猜想、一般函子性与韦伊纲领 II 的算术化，正是这张网尚未织完的下一批结点。
